\documentclass[10pt,a4paper]{amsart}
\usepackage[margin=19mm]{geometry}
\usepackage{amsmath,amssymb,mathtools}
\usepackage[scr=rsfs,cal=euler]{mathalfa}
\usepackage{xfrac}
\usepackage[unicode,hidelinks,bookmarks=false]{hyperref}
\newcommand{\ie}{i.e.\ }
\newcommand{\cf}{cf.\ }
\newcommand{\resp}{respectively\ }
\newcommand{\Subsection}{\subsection}
\providecommand{\nobreakdash}{\nobreak\hbox{-}}
\setlength{\emergencystretch}{2em}
\multlinegap0pt
\usepackage{bm}
\usepackage{bbm}
\usepackage{dsfont}
\usepackage{xspace}
\usepackage{cancel}
\usepackage{mathtools}
\usepackage{amsthm}
\makeatletter
\@ifundefined{corollary}{\newtheorem{corollary}{Corollary}}{}
\@ifundefined{lemma}{\newtheorem{lemma}{Lemma}}{}
\@ifundefined{proposition}{\newtheorem{proposition}{Proposition}}{}
\makeatother
\newcommand{\Bin}[2]{\textbf{Bin}(#1,#2)}
\newcommand{\abs}[1]{\left|#1\right|}
\newcommand{\eps}{\varepsilon}
\renewcommand{\Pr}{\operatorname*{\mathbb{P}}}
\title[Quasipolynomial Trace Reconstruction]{Quasipolynomial Trace Reconstruction}
\author{Arnav Burudgunte, Paul Valiant, Hongao Wang}
\date{Source: arXiv:2607.04073v1 (CC BY 4.0)}
\begin{document}
\maketitle

\noindent\emph{Source and licence.} Quasipolynomial Trace Reconstruction. Version arXiv:2607.04073v1 (CC BY 4.0), distributed under the Creative Commons Attribution 4.0 International licence, as recorded in the arXiv licence metadata for this exact version and shown on the abstract page https://arxiv.org/abs/2607.04073v1. Full rights notice: licenses/arxiv-2607.04073-CC-BY-4.0.txt.\par\smallskip

\noindent\textit{Editorial scope.} Counted unit: the Proposition whose source label is \texttt{prop:induction} in the article (source0.tex lines 1007--1010, source label \texttt{prop:induction}), delivered together with the article's own complete original proof of it (source0.tex lines 1012--1143, one \texttt{proof} environment of 132 lines). No mathematics is written, repaired or re-derived: statement and proof are the article's own text line for line. Required context retained uncounted: lem:binomial-shift (lines 284--287); lem:combined-contrapositive (lines 566--573); cor:shift (lines 592--594); cor:deconvolution (lines 921--924); lem:three-to-one (lines 958--961); lem:bump-pdf-epsilon (lines 1257--1259); lem:cumulative-epsilon (lines 1264--1266). Every definition, display and label of those lines is retained unchanged. Omitted from this package and not redistributed: 1--283, 288--565, 574--591, 595--920, 925--957, 962--1006, 1011--1011, 1144--1256, 1260--1263, 1267--1278. Nothing inside a retained excerpt is omitted or altered: every retained body is delivered exactly as the article's lines are written, so no omission receipt of any kind accompanies this package. Citations: the retained excerpts carry no citation command at all, so no bibliography entry of the article is transcribed and no reference is redistributed. Reference adaptation: none was needed and none was made; every reference inside the delivered unit resolves inside it, so no unresolved reference or citation is produced. Printed slips kept literal: no printed word, symbol, hypothesis or number of the counted statement or of its proof was altered. Layout and class: the article's own class is \texttt{article} with packages including inputenc, amsmath, amsthm, amsfonts, fullpage, amssymb, xcolor, hyperref, bbm, graphicx, algorithm2e, authblk, mathtools, cleveref; the wrapper is the lane's pinned \texttt{amsart} editorial wrapper at 10pt on A4, which restates the theorem-like environments the retained text needs and adds the article's own macro definitions that the retained text needs (\texttt{\textbackslash Bin}, \texttt{\textbackslash Pr}, \texttt{\textbackslash abs}, \texttt{\textbackslash eps}), with extra wrapper packages (bm, bbm, dsfont, xspace, cancel, mathtools). The delivered numbering is the unit's own. Assets: no retained span contains a figure environment, a graphics-inclusion command, a PDF-page inclusion, a table or a TikZ picture; no image file is copied into this package, the package declares no asset dependency and no image file is redistributed with it. Licence: version arXiv:2607.04073v1 of this article is distributed under the Creative Commons Attribution 4.0 International licence, as recorded in the arXiv licence metadata for this exact version and shown on the abstract page https://arxiv.org/abs/2607.04073v1. A full rights notice is delivered at licenses/arxiv-2607.04073-CC-BY-4.0.txt. Characters: every retained excerpt is pure ASCII, so the delivered engine, XeLaTeX with its default Latin Modern text font, reports no missing character.
\begin{lemma}\label{lem:binomial-shift} Let $x \in \{0,1\}^n$. For any statistic $s$ and $0 \leq R \leq n$, we have
\[x_{[1:n]}^{(s)}=\eps+\sum_{j=0}^{R} bin(R,j,p) \cdot x_{[R+1:n]}^{(s-j)}\]
where $0\leq \eps\leq Pr[\Bin{R}{p}\geq \min_i s_i]$.
\end{lemma}

\begin{lemma}\label{lem:combined-contrapositive}
Given an integer $R\geq 250$ and sequences $f,g:\mathbb{Z}\rightarrow\mathbb{R}_{\geq 0}$ supported on $\{-R,\ldots,R\}$ with $\|f\|_1,\|g\|_1\leq \frac{1}{2}$, and a binomial distribution $\Bin{R}{p}$ with $p\in (0,\frac{1}{4}]$, then define the binomially blurred versions of $f,g$ by convolving with the binomial: let $f_b:=f\ast \Bin{R}{p}$ and $g_b:=g\ast \Bin{R}{p}$. Then, if there exists a bound $\tau\in (0,\frac{1}{4}]$ and a location $\ell\leq pR+\sqrt{Rp\log\frac{1}{\tau}}$  such that $|f_b(\ell)-g_b(\ell)|\geq\tau$, and if $\sqrt{\frac{\log\frac{1}{\tau}}{Rp}}\leq \frac{\pi}{8}$ and the left tail of $f-g$ is controlled as $\sum_j |f(j)-g(j)|\cdot e^{-2j\sqrt{(\log\frac{1}{\tau})/(Rp)}}\leq \frac{1}{\tau^c}$ for some constant $c\geq 1$, then: 
\begin{itemize}
    \item[\textup{1)}] Either there exist offsets $\ell_0=0,\ell_1\geq 0,\ell_2\geq 0$ such that $|\sum_j f_b(j+\ell_0)f_b(j+\ell_1)f_b(j+\ell_2)-g_b(j+\ell_0)g_b(j+\ell_1)g_b(j+\ell_2)|\geq 1/poly_c(\frac{1}{\tau},R)$, or the analogous product of two terms ignoring $\ell_2$;
    \item[\textup{2)}] Or there exists $\alpha\in\mathbb{R}$ such that ``$f_b$ is approximately a shifted version of $g_b$'' in the sense that, defining the Fourier transforms $F_b(\xi)=\sum_j f_b(j)\cdot  e^{ij\xi}$ and $G_b$ defined correspondingly in terms of $g_b$, then for all angles $\xi$ that are integer multiples of $\frac{2\pi}{R^2}$ we have $|F_b(\xi)-G_b(\xi)\cdot e^{i\xi \alpha}|\leq \tau^2$.
\end{itemize}

\end{lemma}

\begin{corollary}\label{cor:shift}
Under the assumptions of Lemma~\ref{lem:combined-contrapositive}, if Conclusion 2) applies, then we can further conclude that \[\Big|\sum_{j} (f_b(j)-g_b(j))\cdot e^{ij\frac{2\pi}{R^2}}\Big|\geq \frac{\tau}{R^2}-\tau^2\]
\end{corollary}

\begin{corollary}\label{cor:deconvolution}
Let $f:\mathbb{Z}\rightarrow\mathbb{C}$ supported on an interval $[I_-,I_+]$ with $\sum_{j=I_-}^{I_+} |f(j)| \leq 1$. Further, letting $F:[-\pi,\pi]\rightarrow\mathbb{C}$ be the Fourier transform of $f$, we assume that, for a given $\eps,\delta>0$ we have that $|F(\xi)|\leq\delta$ when $|\xi|\geq \eps$. 
If for $P\in(0,\frac{1}{4}]$ and positive integer $R'$ we have $e^{-0.16PR'+1}\leq\delta\leq e^{-4PR'\eps^2}<1$ then there exists a $g(x)$ that is ``an $r$-local deconvolution of $f$ with $\Bin{R'}{P}$'' in the following sense: for $r=300\sqrt{R'P\lceil\log\frac{1}{\delta}\rceil}$ we have $g$ is supported on $[I_--R'P-r,I_+-R'P+r]$, with the pointwise bound $|g(x)|\leq 11\cdot e^{R'P\eps^2}$, and, $\|g\ast \Bin{R'}{P}-f\|_\infty\leq 9\delta$.
\end{corollary}

\begin{lemma}\label{lem:three-to-one}
Let $x,y$ be two binary strings and let $p\in(0,\frac{1}{4}]$ be a retention probability.
Suppose there is a weight function $w:\mathbb{Z}\rightarrow\mathbb{R}_{\geq 0}$ supported on a positive integer interval $[I_-,I_+]$, having sum $\leq 1$, and for frequencies $|\xi|\geq \eps$ has Fourier transform of magnitude $\leq \delta$, for $\eps\in(0,\frac{1}{3}p]$. Suppose there is a $k$-tuple of distinct positive integers $s$ with $\min_i s_i=1$ and $\max_i s_i=\sigma$ such that, for three offsets $\ell_0=0,\ell_1,\ell_2\geq 0$, defining $f(j):=w(j) x_p^{(s+j)}$ and $g(j):=w(j) y_p^{(s+j)}$, with $|\sum_j f(j)f(j+\ell_1)f(j+\ell_2)-g(j)g(j+\ell_1)g(j+\ell_2)|\geq \tau$ or the corresponding product of two terms omitting $\ell_2$, for some $\tau>0$. Defining $\lambda=3I_++\ell_0+\ell_1+\ell_2$ then: there exists a probability $P\in[p,3p]$, a tuple $S$ of $\leq 3k$ distinct positive integers having minimum value 1 and maximum value $\leq 10(I_++\sigma)$, and there exists a weight function $A:\mathbb{Z}\rightarrow\mathbb{R}_{\geq 0}$ that is supported within $[I_-,\lambda]$, sums to $\leq 1$, and where $A$ has Fourier transform of magnitude $\leq 300 (I_++\sigma)^2\cdot\left(\lambda\delta+\max(\lambda,\frac{8\pi}{3\eps})\cdot e^{-\eps I_-/4}\right)$ at frequencies greater than $9\eps$ such that $|\sum_j A(j)(x_P^{(S+j)}-y_P^{(S+j)})|\geq \frac{\tau-3\cdot 2^{-I_+}}{(10(I_++\sigma))^{3k}}$.
\end{lemma}

\begin{lemma}\label{lem:bump-pdf-epsilon}
Given integer $m >0$ and real number $\eps\in(0,1]$, there exists a function $B^{pdf}:\mathbb{Z}\rightarrow\mathbb{R_{\geq 0}}$ with sum 1 and supported on $(\frac{m}{\eps} - \frac{m}{2}, \frac{m}{\eps} + \frac{m}{2})$ such that for any $|\xi|\geq \eps$, the magnitude of the Fourier transform of $B^{pdf}$ is at most $8\cdot e^{-\eps m/4}$. 
\end{lemma}

\begin{lemma}\label{lem:cumulative-epsilon}
Given integer $m>0$ and real number $\eps\in(0,\pi]$, let $b$ denote the function from Lemma~\ref{lem:bump-pdf-epsilon}, shifted to have support in $\{1,\ldots,m\}$. Then, for an integer interval $I=\{I_-,\ldots,I_+\}$ define the function $B_{I,\eps,m}(j):=\sum_{\ell=I_-}^{I_+-m-1} b(j-\ell)$. Then $B_{I,\eps,m}$ is supported in $I$; equal to 1 on the subinterval that is $m$ smaller, $\{I_-+m,\ldots,I_+-m\}$; has sum $I_+-I_- -m$; takes values in $[0,1]$; and its Fourier transform at angle $|\xi|\geq \eps$ has magnitude at most $\frac{8\pi}{\eps} e^{-m\eps/4}$.
\end{lemma}

\begin{proposition}\label{prop:induction}
Given $n$-bit strings $x,y$ where $d=\min\{i:x_i\neq y_i\}$ is the point of first discrepancy, and a retention probability $p\in(0,\frac{1}{12}]$. For integer $R$ larger than some universal constant, if $\tau$ satisfies $4k\log R\leq \log\frac{1}{\tau}\leq R^{0.8}p^2$ and if there exists a statistic $s$ consisting of $k$ integers in $\{1,\ldots,\sigma\}$ with $\min_i s_i=1$, where $\sigma\leq  \sqrt{Rp\log\frac{1}{\tau}}$  
(``$s$ has order $k$ and span $\sigma$'') such that for some integer $\ell\leq 2Rp+ \sqrt{R p\log\frac{1}{\tau}}$ we have $|x_{[d-2R:n],p}^{(s+\ell)}-y_{[d-2R:n],p}^{(s+\ell)}|\geq \tau$ then: there exists an integer $R'\in[3R^{1.1},R^{2}]$ and a probability $P\in [p,3p]$ and a bound $\tau'\geq 1/poly(1/\tau)$ for which there is an order $\leq 3k$ statistic $S$ of span $\leq  \sqrt{\frac{1}{2}R'P\log\frac{1}{\tau'}}$ and $12k\log R'\leq \log\frac{1}{\tau'}\leq ({R'}/2)^{0.8}P^2$ and a location $\ell'\leq R'P+ \sqrt{\frac{1}{2}R' P\log\frac{1}{\tau'}}$ for which $|x_{[d-R-R':n],P}^{(S+\ell')}-y_{[d-R-R':n],P}^{(S+\ell')}|\geq \tau'$.
\end{proposition}

\begin{proof}
Because we assume that $\log\frac{1}{\tau}\leq R^{0.8}p^2$, we can thus say that for $\emph{any}$ constant $c^*>0$ we have $\log\frac{1}{\tau}\leq c^* R p$ for sufficiently large $R$. We will repeatedly use variants of this bound to show that tail probabilities of the form $e^{-c^*Rp}$ are smaller than $\tau$---or even $\tau$ to any constant power---under the assumptions of the proposition: that $R$ exceeds an appropriately chosen universal constant.

As a first use of this fact, we show that $\ell\in [1.95Rp,2.05Rp]$. The upper bound follows from our assumption that $\ell\leq 2Rp+ \sqrt{R p\log\frac{1}{\tau}}$, since from above, $\log\frac{1}{\tau}$ is smaller than any desired constant multiple of $Rp$. On the other hand since $s+\ell$ has a maximum index $\sigma+\ell$, then a difference in $x_{[d-2R:n],p}^{(s+\ell)}$ versus $y_{[d-2R:n],p}^{(s+\ell)}$ can only happen if some bit at location $\geq d$ in the original strings ends up at location $\leq \sigma+\ell$ in the trace; this happens when $<\sigma+\ell$ of the initial $2R$ bits of the substring $x_{[d-2R:n]}$ are retained in the trace, which happens with probability $\leq \Pr[\Bin{2R}{p}<\sigma+\ell]$, which thus bounds $\tau$. Since by assumption $\sigma\leq  \sqrt{Rp\log\frac{1}{\tau}}$, which is smaller than $0.025Rp$ for large enough $R$, thus if $\ell<1.95Rp$ then we have $\tau\leq  \Pr[\Bin{2R}{p}<1.975Rp]$, which violates our assumptions for large enough $R$, since the right hand side decays exponentially with $Rp$.

We now choose a frequency cutoff $\eps=c' \frac{1}{Rp}\log\frac{1}{\tau}$, for a constant $c'$ to be chosen later (see Equation~\ref{eq:constants}), and use this to choose a smooth function $w:=\frac{1}{Rp} B_{[0.6Rp,1.4Rp],\eps,0.3Rp}$ as defined in Lemma~\ref{lem:cumulative-epsilon}, which from Lemma~\ref{lem:cumulative-epsilon} thus is supported on integers in $[0.6Rp,1.4Rp]$, takes values in the range $[0,\frac{1}{Rp}]$ and is equal to $\frac{1}{Rp}$ on integers in $[0.9Rp,1.1Rp]$, has sum $0.5$, and its Fourier transform at angles $|\xi|\geq \eps$ has magnitude at most $\delta:=\frac{8\pi}{\eps} e^{-0.3 Rp \eps/4}=\frac{8\pi}{\eps} \tau^{0.075 c'}$.

Define the nonnegative sequences \[f(j-Rp):=w(j) x^{(s+j)}_{[d-R:n],p} \quad\textrm{and}\quad g(j-Rp):=w(j) y^{(s+j)}_{[d-R:n],p}\] namely, the expected value of any shift $j$ of the statistic $s$, for a deletion channel of retention probability $p$ starting $R$ to the left of the point of first discrepancy $d$, scaled by the weights $w(j)$, and where we shift the sequences $Rp$ to the left. The shift by $-Rp$ makes the expected location of the point of first discrepancy, in a deletion channel starting at $d-R$, end up at location 0 in $f$ and $g$, so thus we effectively ``center $f,g$ around the expected location of the point of first discrepancy in the trace.''

From the support bounds on $w$, we have that the sequences $f,g$ are supported in $[-0.4Rp,0.4Rp]$. And since statistics take values in $[0,1]$, we have $\|f\|_1,\|g\|_1\leq \sum_j w(j) =0.5$. 

We will apply Lemma~\ref{lem:combined-contrapositive} to the sequences
\begin{equation*}
    f_b := f \ast \Bin{R}{p}, \quad g_b := g \ast \Bin{R}{p}
\end{equation*}
Lemma~\ref{lem:combined-contrapositive} requires certain conditions to hold for the sequences $f,g$ and $f_b,g_b$, which we now verify.

\begin{lemma}
    \label{lem:contrapositive-conditions}
    For $R$ large enough, there exists $\ell \in [1.95RP,2.05Rp]$ and $\tau_{lem} := \frac{\tau}{2Rp}$ such that $\abs{f_b(\ell-Rp) - g_b(\ell-Rp)} \geq \tau_{lem}$, and $\sqrt{\frac{\log \frac{1}{\tau_{lem}}}{Rp}} \leq \frac{\pi}{8}$, and \begin{equation}
        \label{eq:left-tail}
        \sum_j |f(j)-g(j)|\cdot e^{-2j\sqrt{(\log\frac{1}{\tau_{lem}})/(Rp)}}\leq poly\left(\frac{1}{\tau_{lem}}\right)
    \end{equation}
\end{lemma}
\begin{proof}
    We relate $f_b-g_b$ to our assumption that $|x_{[d-2R:n],p}^{(s+\ell)}-y_{[d-2R:n],p}^{(s+\ell)}|\geq \tau$ to set up the application of Lemma~\ref{lem:combined-contrapositive}. From Lemma~\ref{lem:binomial-shift} we have that 
    \begin{equation}
        \label{eq:2R-difference}
        x_{[d-2R:n],p}^{(s+\ell)}-y_{[d-2R:n],p}^{(s+\ell)}=e_x-e_y+\sum_{j=0}^R bin(R,j,p)\left(x_{[d-R:n],p}^{(s+\ell-j)}-y_{[d-R:n],p}^{(s+\ell-j)}\right)
    \end{equation}
    where  $0\leq e_x,e_y\leq \Pr[\Bin{R}{p}\geq \ell+1]$. Since, as proven at the start, $\ell\in [1.95Rp,2.05Rp]$, we have that $|e_x-e_y|$ decays exponentially in $Rp$ and thus is at most $\frac{1}{10}\tau$ for large enough $R$. Next, we will compare the above binomial sum with the very similar expression
    \begin{equation}
        \label{eq:smoothed difference}
        f_b(\ell-Rp)-g_b(\ell-Rp)=\sum_{j=0}^R bin(R,j,p) w(\ell-j) \left(x^{(s+\ell-j)}_{[d-R:n],p}-y^{(s+\ell-j)}_{[d-R:n],p}\right)
    \end{equation}
    When $\ell-j\in[0.9Rp,1.1Rp]$ then $w(\ell-j)=\frac{1}{Rp}$, and thus the summand is exactly $\frac{1}{Rp}$ times the summand in Equation~\ref{eq:2R-difference}. The other case, $\ell-j\notin[0.9Rp,1.1Rp]$, since $\ell\in [1.95Rp,2.05Rp]$, is thus possible only when $j\notin [0.95Rp,1.05Rp]$; and thus we bound the total contribution from this case, using the range bound $w(j)\in [0,\frac{1}{Rp}]$, to conclude that Equation~\ref{eq:smoothed difference} is within $\frac{1}{Rp}\Pr[\Bin{R}{p}\notin [0.95Rp,1.05Rp]]$ of the value of Equation~\ref{eq:2R-difference} without the $e_x-e_y$ terms and when multiplied by $\frac{1}{Rp}$. Thus, from Chernoff bounds we have that $|f_b(\ell-Rp)-g_b(\ell-Rp)|\geq \frac{\tau}{2Rp}$ for $R$ above a large enough constant. 

    We point out that $\sqrt{\frac{\log\frac{1}{\tau_{lem}}}{Rp}}\leq \frac{\pi}{8}$ for sufficiently large $R$, satisfying the second condition.

    Finally, we prove Equation~\ref{eq:left-tail}. As argued above, $|f(j)-g(j)|\leq \frac{1}{Rp}\Pr[\Bin{R}{p}<Rp+j+\sigma]$, which, by Chernoff bounds, is at most $\frac{1}{Rp}\cdot e^{-\frac{(j+\sigma)^2}{3Rp}}$ when $j+\sigma\leq 0$. Thus we can bound the $j\leq -\sigma$ portion of the sum of Equation~\ref{eq:left-tail} by $\frac{1}{Rp}\sum_j e^{-\frac{(j+\sigma)^2}{3Rp}}e^{-2j\sqrt{(\log\frac{1}{\tau_{lem}})/(Rp)}}$. We reexpress the exponential by completing the square, getting the equivalent expression $\frac{1}{Rp}\frac{1}{\tau_{lem}^3}e^{2\sigma\sqrt{\frac{\log\frac{1}{\tau_{lem}}}{Rp}}}\sum_j e^{-\frac{\left(j+\sigma+3\sqrt{Rp\log\frac{1}{\tau_{lem}}}\right)^2}{3Rp}}$. The sum here is bounded by $O(\sqrt{Rp})$, since it decays exponentially fast outside an interval of size $\sqrt{Rp}$ around $j=-\sigma-3\sqrt{Rp\log\frac{1}{\tau_{lem}}}$; meanwhile, the multipliers outside the sum are polynomial in $\frac{1}{\tau_{lem}}$, since, in particular, $\sigma\leq  \sqrt{Rp\log\frac{1}{\tau}}$ by assumption, and thus we bound the exponential factor as $e^{2\sigma\sqrt{\frac{\log\frac{1}{\tau_{lem}}}{Rp}}}\leq \frac{1}{\tau_{lem}^{3 }}$, using the fact that $\log\frac{1}{\tau_{lem}}\leq 2\log\frac{1}{\tau}$ for large enough $R$; thus we get a $poly(\frac{1}{\tau_{lem}})$ bound for the portion of Equation~\ref{eq:left-tail} coming from $j\leq -\sigma$.

    For the remaining portion of the sum, we note that $\sum_j |f(j)-g(j)|\leq 1$, and that, for $j>-\sigma$ we may bound the multiplier $e^{-2j\sqrt{(\log\frac{1}{\tau_{lem}})/(Rp)}}\leq e^{2\sigma\sqrt{(\log\frac{1}{\tau_{lem}})/(Rp)}}\leq \frac{1}{\tau_{lem}^{3 }}$ as above, leading to an overall $poly(\frac{1}{\tau_{lem}})$ bound for Equation~\ref{eq:left-tail}.
\end{proof}

We may now apply Lemma~\ref{lem:combined-contrapositive} to $f,g$. We show that both conclusions of the lemma imply the following general statement about summed statistics. 

\begin{lemma}
    \label{lem:coefficients-helper}
    There exists a sequence $A: \mathbb{Z} \rightarrow \mathbb{C}$ supported on $[0.6Rp,5.8Rp]$ with $\|A\|_1 \leq 1$; there exists a statistic $S$ of min value 1, of order $\leq 3k$ and of span $\leq 15Rp$; and there exists a probability $P \in [p,3p]$ such that 
    \begin{equation}\label{eq:combined-discrepancy-bound}
        \abs{\sum_j A(j) (x^{(S+j)}_{[d-R:n],P} - y^{(S+j)}_{[d-R:n],P}} \geq \tau^{c+1}
    \end{equation}
    for some constant $c \geq 1$. Further, the Fourier transform of $A$ has magnitude $\leq \tau^{0.075c'-1}$ for frequencies $|\xi| \geq 9\eps$. 
\end{lemma}
\begin{proof}
    Having checked the conditions in Lemma~\ref{lem:contrapositive-conditions}, we apply Lemma~\ref{lem:combined-contrapositive} to $f,g$ with threshold $\tau_{lem}=\frac{\tau}{2Rp}$, yielding either Conclusion 1 or Conclusion 2 of Lemma~\ref{lem:combined-contrapositive}. We consider both cases. \\
    
    \noindent \textbf{Conclusion 1}: In this case, there exist offsets $\ell_0=0,\ell_1\geq 0,\ell_2\geq 0$ such that $|\sum_j f_b(j+\ell_0)f_b(j+\ell_1)f_b(j+\ell_2)-g_b(j+\ell_0)g_b(j+\ell_1)g_b(j+\ell_2)|\geq \tau^c$, or the analogous double product---where, since we assume $\frac{1}{\tau}\geq R$, we have absorbed the polynomial dependence on $R$ in the lemma into just a power of $\tau$, for $R$ larger than some absolute constant. We assume $c\geq 1$, and otherwise round $c$ up to 1, preserving the inequality. 

    Recalling that $f_b,g_b$ are the convolution of $f,g$ respectively with $\Bin{R}{p}$, thus a triple product of $f_b$ implicitly consists of an average of triple products of $f$ shifted by a triple of integers $j_0,j_1,j_2\sim \Bin{R}{p}$. Thus, we now remove all terms for which any of $j_0,j_1,j_2\notin [0.9Rp,1.1Rp]$, which changes each of $f_b,g_b$ by an amount exponentially small in $Rp$; for sufficiently large $R$ this change is at most $\frac{1}{2}\tau^c$. By the pigeonhole principle, thus there is a single $j_0,j_1,j_2\in [0.9Rp,1.1Rp]$ whose discrepancy is at least the average, and thus we conclude that, for this $j_0,j_1,j_2\in [0.9Rp,1.1Rp]$ we have \begin{equation}\label{eq:triple-bound}|\sum_j f(j+\ell_0-j_0)f(j+\ell_1-j_1)f(j+\ell_2-j_2)-g(j+\ell_0-j_0)g(j+\ell_1-j_1)g(j+\ell_2-j_2)|\geq \frac{1}{2}\tau^c\end{equation}

    We further point out that $f,g$ both have support $[-0.4Rp,0.4Rp]$, and thus all of $\ell_0-j_0,\ell_1-j_1,\ell_2-j_2$ must be within $0.8Rp$ of each other for Equation~\ref{eq:triple-bound} to be nonzero.

    Next, we apply Lemma~\ref{lem:three-to-one} to the bound of Equation~\ref{eq:triple-bound}, setting $\tau_{lem}:=\frac{\tau^c}{2}$, and using the weight function $w$ supported on the interval from $I_-^{lem}:=0.6Rp$ to $I_+^{lem}:=1.4Rp$; the input offsets are $\ell_0-j_0,\ell_1-j_1,\ell_2-j_2$, though to match the input conditions of the lemma, we subtract off the min of these three from all of them and relabel it as $\ell_0^{lem}=0$, with $\ell_1^{lem},\ell_2^{lem}\in[0,0.8Rp]$ from our above observation that all of $\ell_0-j_0,\ell_1-j_1,\ell_2-j_2$ must be within $0.8Rp$ of each other. Thus in the context of Lemma~\ref{lem:three-to-one} we have $\lambda_{lem}:= 3I_+^{lem}+\ell_1^{lem}+\ell_2^{lem}\leq 5.8Rp$. We check the input condition of the lemma that $\eps\leq\frac{1}{3}p$: since $\eps:=c' \frac{1}{Rp}\log\frac{1}{\tau}$, by our assumption that $\log\frac{1}{\tau}\leq R^{0.8}p^2$ we have $\eps\leq c' R^{-0.2} p$, which is thus $\leq\frac{1}{3}p$ for large enough $R$. We thus invoke the lemma, which yields that: there exists a tuple $S$ of $\leq 3k$ distinct positive integers having minimum value 1 and maximum value $\leq 10(I_+^{lem}+\sigma)\leq 15Rp$ and there exist nonnegative coefficients $A:\mathbb{Z}\rightarrow \mathbb{R}_{\geq 0}$ supported on $[0.6Rp,5.8Rp]$, that sum to $\leq 1$, has discrepancy $|\sum_j A(j)(x_{[d-R:n],P}^{(S+j)}-y_{[d-R:n],P}^{(S+j)})|\geq \frac{\tau^c}{2(15Rp)^{3k}}\geq\tau^{c+1}$, and has Fourier transform that at frequencies $|\xi|\geq 9\eps$ has magnitude  $\leq 300 (I_+^{lem}+\sigma)^2\cdot\left(\lambda_{lem}\delta+\max(\lambda_{lem},\frac{8\pi}{3\eps})\cdot e^{-\eps I_-^{lem}/4}\right)$. Recall from above that $\eps=c' \frac{1}{Rp}\log\frac{1}{\tau}$ and $\delta=\frac{8\pi}{\eps} \tau^{0.075 c'}$. Thus our overall Fourier bound, since $I_-^{lem}=0.6Rp\geq 0.3Rp$, is $\leq \tau^{0.075c'} (Rp)^3$ times a constant, which we can thus bound by $\tau^{0.075c'-1}$ when $R$ exceeds some global constant. \\

    \noindent \textbf{Conclusion 2:} In this case we further apply Corollary~\ref{cor:shift}, which yields that $\big|\sum_{j} (f_b(j)-g_b(j))\cdot e^{ij\frac{2\pi}{R^2}}\big|\geq \frac{\tau_{lem}}{R^2}-\tau_{lem}^2$, which is at least $2\tau^2$ for $R$ larger than some constant since $\tau_{lem}=\frac{\tau}{2Rp}$, and we assumed $\tau\leq \frac{1}{R^4}$ and $p\leq \frac{1}{4}$. 
    Analogously to the Conclusion 1 case, this conclusion is modified by an amount exponentially small in $Rp$ if we limit the binomial convolutions in $f_b,g_b$ to $j'\in[0.9Rp,1.1Rp]$ and we can use pigeonhole to choose a single such $j'$ with at least the average discrepancy. Thus for some $j'\in[0.9Rp,1.1Rp]$ we have $\big|\sum_{j} (f(j-j')-g(j-j'))\cdot e^{ij\frac{2\pi}{R^2}}\big|\geq \tau^2$. However, the shift $j'$ does not affect the magnitude of the discrepancy, so we may drop $j'$ to simply conclude that in the Conclusion 2 case of Lemma~\ref{lem:combined-contrapositive} we have \begin{equation}\label{eq:shift-conclusion}\big|\sum_{j} (f(j)-g(j))\cdot e^{ij\frac{2\pi}{R^2}}\big|\geq \tau^2\end{equation}
    We define $A(j):=w(j) e^{-i(j-Rp)\frac{2\pi}{R^2}}$, so that Equation~\ref{eq:shift-conclusion} says that $|\sum_j A(j)(x_{[d-R:n],p}^{(s+j)}-y_{[d-R:n],p}^{(s+j)})|\geq \tau^2$. Thus this result is analogous to that for the Conclusion 1 case, except we use the original $s,p$ instead of the new $S,P$, and thus the original bounds on $s,p$ are strong enough. We still have strong Fourier bounds on $A$: elementwise multiplying $w(j)$ by $e^{-i(j-Rp)\frac{2\pi}{R^2}}$ simply shifts its Fourier transform by angle $\frac{2\pi}{R^2}$ and multiplies by a phase; thus since since $\frac{2\pi}{R^2}\leq \eps:=c'\frac{1}{Rp}\log\frac{1}{\tau}$ for large enough $R$, we conclude that, in the Conclusion 2 case, $A$ has Fourier transform that at frequencies $|\xi|\geq 2\eps$ has magnitude at most $\frac{8\pi}{\eps} \tau^{0.075 c'} \leq \tau^{0.075 c' - 1}$. 

    Since the desired conclusions hold in either case, this completes the proof.
\end{proof}

We define $R':=c''\frac{R^2 p}{\log\frac{1}{\tau}}$ for a constant $c''$ that we specify now, along with the constant $c'$ used to define $\eps$ at the start of this proof, so that they relate to the constant $c$ of Equation~\ref{eq:triple-bound} induced by Lemma~\ref{lem:combined-contrapositive}. Specifically, we choose $c'$ to be any large enough constant and $c''>0$ to be any constant small enough relative to $c'$ so that
\begin{equation}\label{eq:constants}
0.075c'\geq 3+c\quad \textrm{and}\quad c'c''\leq 0.000004
\end{equation}

\begin{lemma}\label{lem:statistics-deconvolution}
    For large enough $R$, there exists a probability $P \in [p,3p]$, a statistic $S$ of min value 1 and order $\leq 3k$ and span $\leq 15Rp$, a location $j'' \leq R'P + 6.1Rp$ for which
    \begin{equation}\label{eq:final-discrepancy}|x_{[d-R-R':n],P}^{(S+j'')}-y_{[d-R-R':n],P}^{(S+j'')}|\geq \tau^{c_f}\end{equation}
    for any sufficiently large constant $c_f$. 
\end{lemma}
\begin{proof}
    Consider the sequence $A$ given by Lemma~\ref{lem:coefficients-helper} and the associated $P\leq 3p\leq\frac{1}{4}$. We apply Corollary~\ref{cor:deconvolution} to $A(-j)$ to approximately deconvolve the sequence by $\Bin{R'}{P}$. Crucially, we will \emph{flip} $A$ before deconvolving: define the notation $\overleftarrow{A}$ so that $\overleftarrow{A}(j):=A(-j)$. We will denote the output of Corollary~\ref{cor:deconvolution} as $\overleftarrow{A}'$, and thus $A'$ will be the flipped version of the output. The parameters of Corollary~\ref{cor:deconvolution} correspond to the bounds obtained from Lemma~\ref{lem:coefficients-helper}: we let $\delta_{lem}:=\tau^{0.075c'-1}$ and $\eps_{lem}:=9\eps$. 

    We first check the input conditions of Corollary~\ref{cor:deconvolution}. For the condition $e^{-0.16PR'+1}\leq\delta_{lem}$, recall that $\delta_{lem}:=\tau^{0.075c'-1}$ while the left hand side is $e^{-0.16PR'+1}\leq e\cdot e^{-0.16 c''\frac{R^2 p^2}{\log\frac{1}{\tau}}}$; since $\log\frac{1}{\tau}\leq R^{0.8}p$, we bound our expression by $e\cdot e^{-0.16 c''R^{0.4}\log\frac{1}{\tau}}$, which is smaller than $\tau^{0.075c'-1}$ for sufficiently large $R$.

    For the input condition $\delta_{lem}\leq e^{-4PR'\eps_{lem}^2}$, the right hand side is $e^{-4Pc''\frac{R^2 p}{\log\frac{1}{\tau}}(9c' \frac{1}{Rp}\log\frac{1}{\tau})^2}\geq\tau^{3\cdot 4\cdot 9^2c''\cdot {c'}^2}$ which is thus at least $\tau^{0.075c'-1}$ if $972c''\cdot {c'}^2\leq 0.075c'-1$; this equation is implied by the constant conditions in Equation~\ref{eq:constants} since $0.075c'\geq 3$ and thus the right hand side is at least $0.075c'-1\geq 0.05c'$, and thus this condition reduces to $c''c'\leq \frac{0.05}{972}$, which is implied by the second condition of Equation~\ref{eq:constants}.
    
    Thus Corollary~\ref{cor:deconvolution} defines a radius $r_{lem}=300\sqrt{p_{lem}R'\lceil\log\frac{1}{\delta_{lem}}\rceil}$ and shows that there exists a sequence with $|A'(j)|\leq 11\cdot e^{PR'\eps_{lem}^2}$ where the support of $A'$ is the support of $A$ shifted by $R'P$ and expanded by radius $r$; and, crucially, $\overleftarrow{A}'\ast \Bin{R'}{P}$ is within $9\delta_{lem}$ of $\overleftarrow{A}$ pointwise. From our above bounds, the coefficients are bounded as $|A'(j)|\leq 11\cdot e^{p_{lem}R'\eps^2}\leq 11\cdot (\frac{1}{\tau})^{3\cdot 9^2c''\cdot {c'}^2}$. 

    Next, the radius added to the support width by deconvolution is bounded by $r_{lem}=300\sqrt{p_{lem}R'\lceil\log\frac{1}{\delta_{lem}}\rceil}$; since $\lceil\log\frac{1}{\delta_{lem}}\rceil\leq 1+(0.075c'-1)\log\frac{1}{\tau}$ we bound this as $\leq 0.075c'\log\frac{1}{\tau}$ for large enough $R$, and thus $300\sqrt{p_{lem}R'\lceil\log\frac{1}{\delta_{lem}}\rceil}\leq 300\sqrt{3p c''\frac{R^2 p}{\log\frac{1}{\tau}} 0.075c'\log\frac{1}{\tau}}=Rp\cdot 300\sqrt{3\cdot 0.075\cdot c'c''}$. We will show that $r_{lem}\leq 0.3Rp$, which requires $300\sqrt{3\cdot 0.075\cdot c'c''}\leq 0.3$, which is satisfied when $c'c''\leq 0.000004$, which is the second condition of Equation~\ref{eq:constants}. Thus, since Lemma~\ref{lem:coefficients-helper} showed that $A$ is supported within $[0.6Rp,5.8Rp]$, then, shifting this domain by $R'P$ and expanding it by radius $r_{lem}\leq 0.3Rp$ lets us conclude that $A'$ is supported within $[R'P+0.3Rp,R'P+6.1Rp]$.

    Thus since $\overleftarrow{A}'\ast \Bin{R'}{P}$ is within $9\cdot \tau^{0.075c'-1}$ of $\overleftarrow{A}$ pointwise, and the sum of the sizes of their supports is $\leq R'+6R$, we have that $\|\overleftarrow{A}'\ast \Bin{R'}{P}-\overleftarrow{A}\|_1\leq \frac{1}{2}\tau^{0.075c'-2}$ for large enough $R$. Since Equation~\ref{eq:constants} ensures that $0.075c'-2\geq c+1$ then, combining these bounds with Equation~\ref{eq:combined-discrepancy-bound}, for $S$ defined in Lemma~\ref{lem:coefficients-helper} we have that: \begin{equation}\label{eq:binomial-discrepancy}|\sum_{j} (\overleftarrow{A}'\ast \Bin{R'}{P})(-j)\cdot(x_{[d-R:n],P}^{(S+j)}-y_{[d-R:n],P}^{(S+j)})|\geq \frac{1}{2}\tau^{c+1}\end{equation}

    The final step is applying Lemma~\ref{lem:binomial-shift} showing that binomially weighted statistics of substrings $x_{[d-R:n]},y_{[d-R:n]}$ may be reexpressed as statistics of substrings starting $R'$ farther back, thus completing our induction step, ``zooming out'' from scale roughly $R$ to scale roughly $R'$ while preserving a nonnegligible difference in statistics of $x$ versus $y$.

    Explicitly, we expand the left hand side of Equation~\ref{eq:binomial-discrepancy}:
\[|\sum_{j} \sum_{j'} \overleftarrow{A}'(-j-j')\cdot bin(R',j',P)\cdot(x_{[d-R:n],P}^{(S+j)}-y_{[d-R:n],P}^{(S+j)})|\]


    Instead of summing over all $j$, instead define $j''=j+j'$ and equivalently sum over all $j',j''$, where we replace all occurrences of $j$ by the equivalent $j''-j'$:

    \[|\sum_{j''} \sum_{j'} \overleftarrow{A}'(-j'')\cdot bin(R',j',P)\cdot(x_{[d-R:n],P}^{(S+j''-j')}-y_{[d-R:n],P}^{(S+j''-j')})|\]

    We now apply Lemma~\ref{lem:binomial-shift} for each $j''$, to conclude:

    \begin{equation}\label{eq:R'-discrepancy}|\sum_{j''} A'(j'')\cdot (x_{[d-R-R':n],P}^{(S+j'')}-y_{[d-R-R':n],P}^{(S+j'')})|+\sum_{j''} |A'(j'')\Pr[\Bin{R'}{P}\geq 1+j'']|\geq \frac{1}{2}\tau^{c+1}\end{equation}

    

    Since $A'$ is supported on $j''\geq R'P+0.3Rp$, we bound $\Pr[\Bin{R'}{P}\geq 1+j'']$ by standard Chernoff bounds by $e^{-\frac{(0.3Rp)^2}{3R'P}}\leq \tau^{\frac{(0.3)^2}{9c''}}$; thus since the size of the domain of $A'$ is $\leq 6Rp\leq \frac{1}{50\tau}$ for large enough $R$, and $|A'(j'')|\leq 11\cdot (\frac{1}{\tau})^{3\cdot 9^2c''\cdot {c'}^2}$ from above, then, provided that $\frac{(0.3)^2}{9c''}-3\cdot 9^2c''\cdot {c'}^2 \geq c+2$, then we conclude that the second sum of Equation~\ref{eq:R'-discrepancy} is at most $\frac{1}{4}\tau^{c+1}$ and thus the first sum is at least $\frac{1}{4}\tau^{c+1}$. This condition on $c,c',c''$ is implied by Equation~\ref{eq:constants}, since the condition $c'c''\leq 0.000004$ yields $\frac{(0.3)^2}{9c''}\geq 2500c'$ and $-3\cdot 9^2c''\cdot {c'}^2\geq -0.000972c'$, and the condition $c'\geq c+3$ lets us conclude the required inequality $\frac{(0.3)^2}{9c''}-3\cdot 9^2c''\cdot {c'}^2 \geq c+2$. Thus we have shown:

    \[|\sum_{j''} A'(j'')\cdot (x_{[d-R-R':n],P}^{(S+j'')}-y_{[d-R-R':n],P}^{(S+j'')})|\geq \frac{1}{4}\tau^{c+1}\]

    Now we use pigeonhole to extract a single $j''$ for which the discrepancy of the statistics $|x_{[d-R-R':n],P}^{(S+j'')}-y_{[d-R-R':n],P}^{(S+j'')}|$ is large: we divide our bound $\frac{1}{4}\tau^{c+1}$ by the bound $|A'(j'')|\leq 11\cdot (\frac{1}{\tau})^{3\cdot 9^2c''\cdot {c'}^2}$ and our bound on the domain size $6Rp\leq \frac{1}{50\tau}$. Thus, for large enough $R$, there exists $j''\in [R'P+0.3Rp,R'P+6.1Rp]$ for which 
    \begin{equation}|x_{[d-R-R':n],P}^{(S+j'')}-y_{[d-R-R':n],P}^{(S+j'')}|\geq \tau^{c+2+3\cdot 9^2c''\cdot {c'}^2}\end{equation}
    Choosing $c_f$ to be at least $c+2+3\cdot 9^2c''\cdot {c'}^2$ yields the lemma.
\end{proof}



It remains only to prove bounds on the statistic $S$ and the values $R',\tau'$ as desired by the proposition statement. From Lemma~\ref{lem:statistics-deconvolution}, we have constructed $S$ as having min value 1 and span $\leq 15Rp$ and order $\leq 3k$.



Define $\tau':=\tau^{c_f}$, where we will choose the constant $c_f$ to be at least $c+2+3\cdot 9^2c''\cdot {c'}^2$ so that Equation~\ref{eq:final-discrepancy} yields a $\tau^{c_f}$ discrepancy bound.

To prove the remaining bounds, we first note that $R':=c''\frac{R^2 p}{\log\frac{1}{\tau}}\in [3R^{1.1},R^2]$ for sufficiently large $R$. Then we check that our bound $15Rp$ on the span of $S$ is at most the requirement $ \sqrt{\frac{1}{2}R'P\log\frac{1}{\tau'}}$; this expression is at least $Rp \cdot  \sqrt{\frac{1}{2}c'' c_f}$, which is thus at least $15Rp$ for large enough constant $c_f$. The location $\ell'$ in the conclusion is defined to be $j''$ which we bounded as $\leq R'P+6.1Rp$, and thus is also bounded by $R'P+ \sqrt{\frac{1}{2}R'P\log\frac{1}{\tau'}}$, for large enough constant $c_f$.

We next show $12k\log R'\leq \log\frac{1}{\tau'}\leq ({R'}/2)^{0.8}P^2$. For the first inequality, recall our assumption that $4k\log R\leq \log\frac{1}{\tau}$; since $R'\leq R^2$, and $\log\frac{1}{\tau'}=c_f\log\frac{1}{\tau}$, we have that this inequality is satisfied when $c_f\geq 6$. For the second inequality, recall our assumption $\log\frac{1}{\tau}\leq {R}^{0.8}p^2$; since $P\geq p$ and since $R'\geq R^{1.1}$ from above, we have that $({R'}/2)^{0.8}P^2$ exceeds any constant $c_f$ times ${R}^{0.8}p^2$, for large enough $R$. We have thus concluded all the desired bounds.
\end{proof}

\end{document}
